\subsection{Closed factors and arbitrary integrable densities}
\label{subsec:closed-factor-equations}

The preceding examples use an explicit potential on projective space. On a closed K\"ahler factor, Stokes' theorem instead bounds the non-pluripolar mass, and this bound makes every threshold subsolution on the factor exact. This gives equations for arbitrary integrable densities on that factor, with no assumption that they are bounded below.

In this subsection $(Y,\kappa)$ is again an arbitrary connected compact K\"ahler manifold with $\int_Y\kappa^k=1$, with the notation $X$, $N$, $H$, $C$, $V_0$, $\mu$ and $R$ introduced at the beginning of the section, and we put $c_0=CV_0$.

\Needspace{12\baselineskip}
\begin{theorem}\label{thm:closed-factor-integrable-equation}\index{Monge--Ampere equation@Monge--Amp\`ere equation!closed factor}\index{Normalization scalar!closed factor}
For every nonzero nonnegative $g\in L^1(Y,\kappa^k)$, with $M_g=\int_Yg\kappa^k$, one has
\[
 b_H(Rg)=\frac{c_0}{M_g}.
\]
There is $v\in\PSH(Y,\kappa)$ with $\sup_Yv=0$ and $\int_Y\mathrm{MA}^{\mathrm{ext}}_\kappa(v)=1$ whose pullback $u$ satisfies
\[
 \mathrm{MA}^{\mathrm{ext}}_H(u)
       =\frac{c_0}{M_g}Rg\,\mathrm dV,\qquad
 \int_X\mathrm{MA}^{\mathrm{ext}}_H(u)=c_0.
\]
\end{theorem}
\begin{proof}
The formula for $b_H(Rg)$ is part (2) of \cref{lem:split-product-normalization}. By \cref{prop:individual-normalizer-lp}, applied on $Y$ with the K\"ahler form $\kappa$ and the reference volume $\kappa^k$, $b_\kappa(g)=1/M_g$. Thus \cref{thm:individual-density-subsolution} on $Y$, followed by subtraction of a constant, gives $v\in\PSH(Y,\kappa)$ with supremum zero and
\[
 \mathrm{MA}^{\mathrm{ext}}_\kappa(v)\geq(g/M_g)\kappa^k.
\]
For $v_j=\max(v,-j)$, Stokes' theorem for the closed form $\kappa$ gives $\int_Y\mathrm{MA}_\kappa(v_j)=1$. Since $\mathrm{MA}^{\mathrm{ext}}_\kappa(v)$ is the increasing limit of the measures $\mathbf1_{\{v>-j\}}\mathrm{MA}_\kappa(v_j)$ in \eqref{eq:extended-ma}, it follows that $\int_Y\mathrm{MA}^{\mathrm{ext}}_\kappa(v)\leq1$. The prescribed lower measure has mass one, so equality holds both in mass and in the measure inequality. Part (1) of \cref{lem:split-product-normalization} now gives
\[
 \mathrm{MA}^{\mathrm{ext}}_H(u)
   =C h_0^m\wedge\mathrm{MA}^{\mathrm{ext}}_\kappa(v)
   =\frac{c_0}{M_g}Rg\,\mathrm dV,
\]
a measure of mass $c_0$.
\end{proof}

When $v_-(h_0)=0$, the background $H$ has zero lower mass and is not closed, by part (3) of \cref{lem:split-product-normalization}. For $Y=\mathbb P^k$ and $g=q_\alpha$, the potential $u_\alpha$ of \cref{thm:zero-lower-mass-product-equation} is such a solution. For the elliptic fivefold of \cref{cor:elliptic-fivefold-exact-threshold}, $c_0=240\pi$; thus, for every probability density $g$ on $\mathbb P^2$, the density $Rg$ has threshold $240\pi$ and admits a threshold solution of finite mass. The theorem concerns factor densities; it does not give an exact equation for an arbitrary density on $X$ or identify limits of normalized exponential solutions.
